\begin{abstract}
\noindent
This report investigates whether single-cell RNA-seq, combined with foundation-model embeddings, can
predict pharmacological response, as a first step toward a reusable pan-cancer single-cell foundation
model. On the highest-confidence intersection of the SCP542 single-cell atlas and the CTRPv2
drug-response database (\NLines{} cell lines, \NDrugsModelled{} compounds), we train a shared-trunk multi-task
regressor on individual cells and evaluate it under a leave-cell-line-out protocol, comparing scGPT
embeddings against a matched principal component analysis (PCA) baseline. The guiding hypothesis is that
the foundation model provides a denoised, tissue-agnostic representation of cellular state.

\revision{\textbf{Quantitative results are withheld in this version.} The response target, the compound
panel and the preprocessing of both representations were each revised after the previous results were
produced, and no measurement has yet been repeated on the corrected pipeline; every number is therefore
withdrawn rather than restated with a caveat (Section~\ref{sec:results}). What the report contains at
this stage is the design and its justification: the provenance and harmonization of the two data sources,
the derivation of the response target, the construction of the compound panel from external evidence, and
the evaluation protocol.}

The methodological argument does not depend on the withdrawn measurements and is the contribution this
version claims. A representation comparison of this kind is only informative once three conditions hold,
and none held at the outset of this work: the training objective must not let a minority of compounds
with large response variance dominate a shared multi-task loss; the evaluation must be scored against a
control that already absorbs the cell-line effect, since a per-drug-mean null is too weak a bar to
separate a representation from an average; and the compounds must be selected without reference to the
response values the model is trained on, or the panel is enriched for exactly the drugs that flatter it.
Each is a route by which a conventional setup reports a representation advantage that is not one --- the
mechanism the DrEval study \cite{dreval} identifies across the published literature. Establishing those
conditions, rather than the comparison they are prerequisites for, is what this stage of the work
delivers.
\end{abstract}
